\scanpage{127}
\mspage{\handwritten{3}}
\figure{\begin{tikzpicture}[x=0.015cm,y=0.015cm]
\draw[line width=0.4pt] (16,56) .. controls (3,54) and (0,36) .. (8,24) .. controls (17,8) and (44,3) .. (79,13);
\end{tikzpicture}}

De ces deux propriétés résulte formellement qu'on peut trouver un systême projectif \inserted{\handwritten{filtrant}} $(G_i)_{i\handwritten{\in}I}$ dans $\underline{G}$, avec des morphismes $G_i\handwritten{\to}G_j$ qui sont des épimorphismes dans $\unsure{\underline{\mathbb{C}}|\underline{\mathbb{G}}}$, ``essentiellement unique'' dans un sens facile à préciser, tel que \deleted{\unsure{le} foncteur $Z$ soit \unsure{isomorphe}} l'on ait un isomorphisme fonctoriel
\[
Z(G)=\underset{\handwritten{\longrightarrow}}{\lim}\operatorname{Hom}_{\handwritten{\underline{G}}}(G_i,G)
\]

C'est ce systême projectif (pris modulo une équivalence qui intuiti\unsure{v}ment signifie qu'on a passé à la limite projective des $G_i$) qui peut être noté $\handwritten{\pi}^{\underline{G}}_1(S,a)$ et joue le rôle d'un groupe fondamental de $S$ en $a$ (relativement à la catégorìe $\underline{G}$ de groupes, munie du foncteur $F$). Dans le cas \deleted{\unsure{pris}} d'un corps de base $k$, $\underline{G}$ étant la catégorie des groupes algébriques finis sur $k$, on po\uline{urra écrir}e $\handwritten{\pi_1}(S/k,a)$,\margin{\handwritten{à tort}} c'est le groupe proalgébrique (mais avec partie infinitésimale) fondamental du $k$-schéma $S$. Lorsque $k$ est parfait, la décomposition d'un groupe algébrique fini en partie infinitésimale et partie réduite montre que le groupe fondamental est produit semi-direct de sa partie réduite ou séparable \deleted{\unsure{définie}} (qui correspond à un groupe compact discontinu ordinaire\handwritten{,} surlequel le groupe fondamental ordinaire de $k$ - i.e. le groupe de Galois de $\bar{k}$ sur $k$ - opère\inserted{\handwritten{lorsque $a$ est \unsure{rationnel} sur $k$}}), par sa partie infinitésimale, \deleted{\unsure{qui sera un p-groupe fi|qui sera un pro-groupe fi}} qui elle ne dépend d'ailleurs plus du choix du point base $a$ de $S$. Noter que la partie séparable du groupe fondamental peut se réconstituer facilement à l'aide du groupe fondamental ordinaire de $S$ et de son homomorphisme naturel dans $\unsure{\deleted{\mathrm{g}}\mathrm{G}|\deleted{\mathrm{G}}\mathrm{g}}\mathrm{al}(\bar{k}/k)$, \inserted{\handwritten{mais}} il est ``plus grand'' que le groupe fondamental ordinaire, parce\inserted{\handwritten{\,}}qu\deleted{e}\inserted{\handwritten{'}} il correspond à la classification de revêtements principaux de groupe structural non seulement un groupe fini ordinaire, mais un groupe \deleted{\unsure{fini sur lequel}} fini et séparable sur $k$ (i.e. un groupe fini ordinaire sur lequel on fait opérer $\operatorname{Gal}(\bar{k}/k)$ de façon non nécessairement triviale). \deleted{\unsure{Il est}}

J'avoue que si $k$ n'est pas algébriquement clos, le groupe fondamental ci-dessus n'est \deleted{\unsure{peut-être}} \inserted{\handwritten{sans-doute}} pas le bon; on devrait sans doute prendre le groupe fondamental de $S\unsure{\otimes|\times}_{\handwritten{k}}\bar{k}$, qui est un groupe proalgébrique défini sur $\bar{k}$, noter qu'il est muni d'une donnée de descente de $\bar{k}$ à $k$ et est par suite défini en réalité sur $k$; \deleted{avec cette définition, si $S=\operatorname{Spec}(k)$, la \deleted{\unsure{groupe}} partie séparable du groupe fondamental de $S$ (relativement à la ponctuation de $S$ définie par une \deleted{extension \unsure{algé}} clôture algébrique $\bar{k}$ de $k$) n'est autre que}
% archive sheet number: 126
